\chapter{PHIL Ledger --- Philosophy Tags}
\label{chap:PHIL_Ledger}

\section{Purpose}

PHIL tags encode philosophical interpretation: ethics, metaphysics, logic,
aesthetics, and modes of reasoning.

\section{Scope}

Describes:
\begin{itemize}
    \item ethical principles,
    \item metaphysical judgments,
    \item reasoning styles,
    \item phenomenological interpretation,
    \item virtue or value commitments.
\end{itemize}

\section{Schema}

\[
T_{\text{PHIL}} = \langle key,\, gloss,\, bindings,\, constraints \rangle
\]

Examples:
\begin{itemize}
    \item \texttt{phil:ethics}
    \item \texttt{phil:logic}
    \item \texttt{phil:phenomenology}
    \item \texttt{phil:virtue}
    \item \texttt{phil:metaphysics}
\end{itemize}

\section{Class Binding Notes}

\begin{itemize}
    \item P-class dominant for value structures.
    \item O-class for philosophical narrative lenses.
    \item S-class for formal logic structures.
\end{itemize}

\section{Constraints}

\begin{itemize}
    \item No prescriptive moral claims without HITL.
    \item Must not contradict governing charters.
    \item PHIL tags affecting Kernel require P-class review.
\end{itemize}
